Automated discovery often alternates proposals with experiments. Can the
consequences of earlier experiments improve the first algorithm proposed for
a new problem, before that algorithm is tested? We define
\emph{InnovationZero} as a data-to-algorithm protocol: a frozen generator
receives permitted training data, commits one reusable learning procedure,
and receives no candidate-specific execution feedback before commitment.
Independent compilation, fitting and evaluation follow. An audit of JopeDime
distinguishes executable construction, catalog-constrained neural learning,
four frontier-labelled regression artifacts and exploratory benchmark
evidence. These establish components of the research program, not its central
generalization claim. In particular, a decoder restricted to a finite catalog
cannot establish synthesis outside that catalog. We specify a prospective,
family-separated comparison with selection, retrieval, conditional generation
and strong fixed learners; failures remain in the denominator. Structural,
behavioral and intervention-based novelty tests distinguish a useful new
procedure from a new source file. Because decisions are judged by quantiles,
weighted deviances and downside risk rather than squared error, the objective
is treated as part of the input, with implemented objective-aware selection,
calibration and program ranking as baselines. Finally, compatible token
accounting and symmetric cost amortization separate retained output from
discovery efficiency. The resulting framework asks how much experimental work can move
into training without confusing a promising first proposal with verified
new knowledge.
